\chapter{Conclusioni e Sviluppi Futuri}
\label{chap:conclusioni}

\section{Sintesi dei Contributi Scientifici e Ingegneristici}
\label{sec:sintesi_contributi}

Il presente lavoro di tesi e progetto d'esame ha affrontato in modo sistematico i limiti strutturali insiti nelle architetture convenzionali di \textit{Retrieval-Augmented Generation} (Naïve RAG) applicate all'Information Retrieval complesso, proponendo un framework modulare basato sul paradigma emergente dell'\textbf{Agentic AI} \cite{zhang2024ai, wu2025agentic}. 

L'obiettivo primario era la progettazione e la validazione empirica di un agente autonomo capace di colmare il divario informativo (\textit{information gap}) che caratterizza le interrogazioni composite multi-hop, le quali non possono essere risolte mediante una scansione documentale a singolo passaggio.

I contributi sviluppati nel corso della ricerca toccano diversi aspetti metodologici e ingegneristici. Anzitutto, abbiamo superato la rigidità delle pipeline sequenziali modellando il sistema come una macchina a stati finiti con memoria transazionale condivisa (\texttt{AgentState}) realizzata tramite LangGraph; l'adozione di guardie numeriche sui cicli e di limiti di ricorsione a runtime assicura una convergenza deterministica in tempo polinomiale, eliminando il rischio di loop infiniti. In secondo luogo, attraverso il nodo supervisore Validator abbiamo implementato i principi di Adaptive-RAG, classificando la complessità dell'intento dell'utente al turno iniziale per evitare costose scomposizioni su quesiti semplici e attivando il raffinamento solo a fronte di domande comparative o a ponte.

A questo si affianca l'integrazione del Least-to-Most Prompting nel nodo Refiner, che permette di scomporre richieste complesse in sotto-query atomiche auto-consistenti e topologicamente ordinate, preservando i nomi delle entità ed evitando la deriva semantica. Sul fronte dell'esecuzione, il nodo Planner adotta il ciclo ReAct coordinando dinamicamente la Knowledge Base locale basata su BM25Okapi (indicizzata su SQLite per risposte ad altissima velocità) e la ricerca live sul Web per colmare tempestivamente qualsiasi lacuna documentale. La fedeltà delle risposte è garantita da un doppio livello di verifica ispirato a Self-RAG, sia tramite sintesi documentata nel Planner, sia mediante il controllo dei token di critica $\text{IsRel}$ e $\text{IsSup}$ nel Validator, con una verifica stringente sulla completezza polare per i quesiti chiusi (sì/no). Infine, lo sviluppo di una dashboard interattiva in Streamlit offre una trasparenza forense completa, consentendo l'ispezione delle transizioni di stato, l'analisi dei singoli hop di pensiero e il confronto differenziale automatico tra modalità di retrieval.

\section{Riflessioni sulle Evidenze Sperimentali e Implicazioni per il Knowledge Management}
\label{sec:riflessioni_implicazioni}

La validazione quantitativa condotta sul subset di validazione di **HotpotQA** (\textit{fullwiki} split) ha fornito riscontri empirici di chiaro valore, dimostrando la superiorità sistematica del Sistema Integrale multi-agente rispetto alla Baseline locale: la Grounded Accuracy cresce del $+16,66\%$ (salendo dal $26,67\%$ al $43,33\%$), il Macro F1-Score a livello di token passa dal $23,67\%$ al $40,93\%$ ($+17,26\%$), e l'Exact Match fa registrare una vera e propria triplicazione, passando dal $13,33\%$ al $40,00\%$ ($+26,67\%$). Al contempo, il numero medio di hop di pianificazione scende da $2,13$ a $1,60$ per interrogazione, confermando che una query decomposta e ripulita massimizza l'efficienza ed evita ricerche dispersive.

Sul piano concettuale, questi risultati segnano il definitivo superamento del modello tradizionale di ricerca basato su una corrispondenza statica tra termini di query e indici invertiti. Nei sistemi di nuova generazione, l'Information Retrieval evolve verso un processo cognitivo iterativo. L'utente non è più costretto ad articolare manualmente query Booleane complesse o a eseguire mentalmente il multi-hop collegando decine di schede nel browser: l'agente si fa carico della formulazione strategica del piano di ricerca. Inoltre, nei contesti aziendali mission-critical, una risposta generativa opaca è inaccettabile; il nostro framework assicura una tracciabilità forense integrale, permettendo di risalire per ogni affermazione al nodo deliberativo, all'hop ReAct, allo strumento impiegato e all'esatto snippet con URL di provenienza.

\section{Analisi Critica delle Limitazioni del Framework}
\label{sec:limitazioni}

Un'analisi ingegneristica trasparente impone di evidenziare alcune limitazioni operative del sistema attuale. 

La prima riguarda la latenza computazionale: l'attraversamento di un grafo multi-agente comporta molteplici invocazioni successive al modello linguistico (per la classificazione iniziale, il raffinamento, i singoli hop ReAct e la sintesi finale). Sebbene il modello locale quantizzato Mistral-Nemo operi con tempi contenuti su GPU, la latenza aggregata per query complesse si attesta tra i 10 e i 20 secondi, risultando superiore rispetto al recupero millisecondo di un indice lessicale puro.

La seconda limitazione risiede nella sensibilità alla qualità e alla frammentazione dei risultati Web: quando il Planner interroga sorgenti aperte, il recupero dipende dai payload restituiti dalle API. Pagine protette da paywall, frammenti testuali interrotti o tabelle complesse possono indurre l'agente a un'interpretazione parziale del dato. A ciò si aggiungono difficoltà su granularità numeriche e relazioni temporali, dove il modello talvolta confonde capienze complessive con posti a sedere o fatica nei confronti di età se manca la data di nascita di uno dei soggetti. Infine, l'accumulo progressivo di testi all'interno del contesto condiviso può avvicinarsi al limite operativo della finestra di contesto per ricerche con molti hop, indicando l'opportunità di introdurre moduli di compattazione periodica.

\section{Direzioni Future e Prospettive di Ricerca}
\label{sec:sviluppi_futuri}

Le limitazioni riscontrate aprono interessanti linee di sviluppo per la prosecuzione della ricerca. Una prima estensione consiste nell'integrazione di basi di conoscenza semantiche e Knowledge Graphs: l'aggiunta di strumenti per l'interrogazione SPARQL o Cypher su grafi strutturati come Wikidata consentirebbe al Planner di risolvere sotto-query relazionali mediante deduzioni logiche formali, azzerando le ambiguità del testo libero. 

Una seconda linea riguarda la sperimentazione di architetture di dibattito multi-agente (\textit{Multi-Agent Debate}): al posto di un singolo supervisore, è possibile configurare un comitato di agenti critici specializzati (un fact-checker, un controllore temporale e un verificatore di consistenza quantitativa) che dibattono collegialmente prima di convalidare la risposta finale. Infine, un traguardo di grande rilievo è l'ottimizzazione end-to-end con Reinforcement Learning: l'addestramento diretto dei pesi dell'agente tramite algoritmi di Policy Gradient (PPO o DPO), impiegando funzioni di ricompensa basate sul grounding e penalità sui costi computazionali per minimizzare sia le allucinazioni sia il numero di hop necessari.

\section{Possibili Estensioni per la Tesi di Laurea Magistrale}
\label{sec:estensioni_tesi}

In piena aderenza con le indicazioni metodologiche formulate dai docenti del corso di Information Retrieval (sezione dedicata alle estensioni di tesi nel documento di specifica), l'architettura agentica sviluppata in questo progetto costituisce una base concettuale solida che può essere estesa verso quattro progetti di **Tesi di Laurea Magistrale** di ampio respiro scientifico.

\subsection{1. Decomposizione Agentica della Query per l'Analisi Spazio-Temporale}
Questa linea di ricerca prevede la progettazione di un planner agentico gerarchico capace di decomporre autonomamente interrogazioni complesse basate su serie storiche di Osservazione della Terra (Earth Observation) in una sequenza ordinata di sotto-task specializzati. Si pensi a richieste articolate come identificare le aree a forte deforestazione nella foresta amazzonica e recuperare le immagini satellitari multispettrali per evidenziare le variazioni stagionali nell'ultimo quinquennio. Sul piano metodologico, il lavoro richiede l'implementazione di un modulo di pianificazione con decomposizione gerarchica dello spazio e del tempo, l'integrazione di interfacce che interrogano cataloghi di immagini satellitari (Copernicus Open Access Hub, Sentinel-2, Landsat) e il coordinamento tramite un gestore di strumenti che alterna algoritmi di Content-Based Image Retrieval (CBIR), change detection e filtri geospaziali sui metadati.

\subsection{2. Tool-Augmented Agentic Retrieval con Geo-Analytics}
Una seconda direzione si focalizza sull'incorporamento di un livello di orchestrazione agentica capace di selezionare e invocare dinamicamente tool computazionali geospaziali avanzati all'interno della pipeline di recupero \cite{lee2025multiagent}. L'approccio metodologico prevede la costruzione di un catalogo di strumenti specialistici che mappa le chiamate dell'agente verso funzioni analitiche concrete: calcolatori di indici di vegetazione (NDVI, NDWI), elaboratori di modelli digitali di elevazione (DEM) e script di geoprocessing in ambiente QGIS o GDAL. L'infrastruttura si connette in tempo reale con API cloud massive per l'elaborazione su scala planetaria, quali Google Earth Engine (GEE), sfruttando prompt ReAct strutturati per correlare grandezze ambientali complesse, come lo stress della biomassa in rapporto a gradienti altimetrici ed eventi siccitosi.

\subsection{3. Fact-Checking Auto-Riflessivo nel Retrieval Multimodale}
La terza proposta estende il nodo Validator verso un agente di verifica multimodale che esegue un controllo incrociato continuo tra immagini satellitari ad alta risoluzione, metadati cartografici e fonti documentali scientifiche prima di licenziare la risposta. La metodologia si fonda su routine di verifica incrociata multi-sorgente per contrastare allucinazioni visive, confrontando ad esempio la classificazione del suolo derivata da RAG multimodale con registri catastali e pubblicazioni idrologiche. Il modello stima un punteggio di confidenza bayesiano che attiva interrogazioni di approfondimento in caso di elevata incertezza, producendo al termine un sommario esplicativo di fact-checking corredato da annotazioni sull'origine delle evidenze e intervalli di confidenza.

\subsection{4. Sistemi di AI Agentica per la Rilevazione e Mitigazione di Vulnerabilità Software}
La quarta estensione, condotta in collaborazione con i laboratori di ricerca di Ingegneria del Software del DIETI, punta allo sviluppo di un agente autonomo specializzato nell'audit continuo di artefatti software, nell'analisi semantica del codice e nella formulazione di contromisure contro falle di sicurezza. L'approccio adotta modelli linguistici orientati al codice, capaci di interpretare la semantica di librerie legacy e moduli di autenticazione attraverso grafi di flusso di controllo e alberi sintattici astratti (AST). L'agente effettua simulazioni di catene di exploit guidate da esecuzione simbolica e algoritmi di reinforcement learning, interrogando dinamicamente basi di dati pubbliche su vulnerabilità note (CVE, NVD, CWE) e alimentando un modulo di sintesi automatica di patch correttive con validazione formale dell'assenza di regressioni funzionali.

L'integrazione di capacità decisionali, strumenti computazionali e auto-valutazione trasforma l'Information Retrieval da una tecnologia di semplice indicizzazione a un **partner cognitivo attivo**, capace di navigare la complessità del patrimonio informativo moderno con rigore, efficienza e trasparenza scientifica.
